\documentclass[conference]{IEEEtran}
\IEEEoverridecommandlockouts
\usepackage{cite}
\usepackage{amsmath,amssymb,amsfonts}
\usepackage{algorithmic}
\usepackage{graphicx}
\usepackage{textcomp}
\usepackage{xcolor}
\usepackage{multirow}
\usepackage{array}
\usepackage{booktabs}
\usepackage{adjustbox}

\def\BibTeX{{\rm B\kern-.05em{\sc i\kern-.025em b}\kern-.08em
    T\kern-.1667em\lower.7ex\hbox{E}\kern-.125emX}}

\begin{document}

\title{Stylmorph: A Virtual Try-On Platform with 2D-to-3D Workflow for Enhanced E-Commerce Experience}

\author{
\IEEEauthorblockN{Devi Nandana, Neeraj S, Sanjay Kumar P M, Vinaya Krishnan U, Mrs. Jincy Showri}
\IEEEauthorblockA{\textit{Department of Computer Science and Engineering} \\
\textit{College of Engineering, Cherthala} \\
Alappuzha, Kerala, India \\
\{devinandanas22, neerajleo19, sanjaykumarpm006, vinayakrishnan107\}@gmail.com}
}

\maketitle

\begin{abstract}
Stylmorph is a comprehensive virtual try-on platform designed to address the lack of personalization and fit visualization in modern e-commerce. The system integrates a robust two-stage processing workflow combining high-speed 2D virtual try-on with detailed 3D avatar generation. Initially, high-fidelity 2D try-on using the FASHN-VTON framework synthesizes selected garments on the user's uploaded photograph after executing on an external Google Colaboratory environment due to infrastructure constraints. Subsequently, the system generates a personalized 3D avatar from the same reference image by leveraging the PIFuHD monocular 3D reconstruction model. By providing detailed 3D avatars for comprehensive 360-degree garment fitting analysis, Stylmorph significantly mitigates size and fit discrepancies. This interactive and personalized shopping experience aims to reduce the high product return rates burdening the retail sector and substantially augment consumer confidence in online apparel purchases.
\end{abstract}

\begin{IEEEkeywords}
Virtual Try-On, 3D Avatars, Monocular 3D Reconstruction, PIFuHD, FASHN-VTON, E-commerce, Computer Vision, Deep Learning, WebGL
\end{IEEEkeywords}

\section{Introduction}

Online fashion shopping has grown rapidly in recent years. However, many customers still face difficulty understanding how clothes will fit their body. Most e-commerce platforms rely on 2D images and size charts, which often do not represent real body shapes accurately. As a result, customers hesitate to buy products, and return rates in online fashion retail remain high \cite{b7}.

Earlier virtual try-on systems required users to manually enter detailed body measurements to generate basic 3D models \cite{b12}. Although useful, these systems lacked realism and personalization. To address this issue, Stylmorph introduces a two-stage approach that combines 2D visualization with detailed 3D avatar generation. Users can upload a single full-body image, and the system first applies 2D virtual try-on using FASHN-VTON \cite{b23} to visualize clothing on the user's photograph. Subsequently, the same image is used to generate a personalized 3D avatar \cite{b4} for detailed fitting analysis.

The platform uses deep learning models such as PIFuHD \cite{b1} and parametric body models \cite{b3} to reconstruct 3D human geometry from monocular images, without requiring multi-camera setups or depth scanners. The frontend is built with React.js and Vite, while the backend combines Node.js and Python for ML execution. Generated 2D try-on results and 3D meshes are stored via Cloudinary and rendered in the browser using Three.js \cite{b8} for interactive 3D visualization.

\subsection{Objectives}
The key objectives of Stylmorph are: (1) to provide instant 2D try-on results for rapid garment browsing; (2) to generate detailed, topologically accurate 3D avatars from a single photo without specialized hardware; and (3) to enable real-time, 360-degree interactive inspection of how garments fit the user's body within a standard web browser.

\subsection{Challenges}
The core challenges include: (1) the high computational overhead of running deep-learning inference models within a real-time web context, necessitating an external Google Colab environment for FASHN-VTON processing due to infrastructure limitations; (2) maintaining accurate 3D body topology from a single 2D image without multi-view depth inputs; and (3) efficiently streaming large binary \texttt{.obj} mesh files from cloud storage to a WebGL client renderer.

\subsection{Paper Organization}
The remainder of this paper is organized as follows. Section II reviews related work in virtual try-on and 3D human reconstruction. Section III presents the proposed system design and methodology, covering the overall workflow, architecture, and individual modules. Section IV describes the experimental setup and implementation environment. Section V presents and discusses the obtained results, including performance metrics, qualitative outputs, and a comparative analysis. Finally, Section VI concludes the paper and outlines directions for future work.

\section{Related Work}

Virtual try-on research has progressed significantly over the last decade due to advancements in computer vision and 3D modeling. Early systems relied on 2D overlays and simple body measurement tools \cite{b12}. While functional, these systems could not provide realistic fitting visualization.

Recent advances in image-based virtual try-on have produced sophisticated frameworks for 2D garment transfer. FASHN-VTON \cite{b23} introduced a novel approach for high-fidelity virtual clothing transfer that preserves garment details, textures, and folds while adapting to the target person's pose and body shape. This builds upon earlier work in 2D virtual try-on such as VITON \cite{b24} and CP-VTON \cite{b25}, which established the foundation for image-based clothing transfer.

The introduction of parametric body models such as SMPL allowed researchers to represent human body shapes more accurately. However, these models typically generate simplified bodies without detailed clothing or textures \cite{b3}. Subsequent work by Kanazawa et al. \cite{b14} and Pavlakos et al. \cite{b15} enabled end-to-end recovery of human shape and pose from single images.

Later, monocular 3D reconstruction methods improved realism. PIFuHD \cite{b1} enabled high-resolution 3D human digitization from a single image, preserving details such as clothing folds and hair. Bogo et al. \cite{b16} developed methods for automatic estimation of 3D human pose and shape, while Lassner et al. \cite{b17} introduced generative models of people in clothing.

Research by Yang et al. \cite{b18} on universal human parsing and Pons-Moll et al. \cite{b19} on dynamic human shape in motion further advanced the field. Early work by Guan et al. \cite{b20} on dressing virtual humans laid the foundation for modern virtual try-on systems.

Although many research prototypes demonstrate strong technical performance, few systems provide integrated 2D and 3D virtual try-on in a web-based deployment. Stylmorph focuses on bridging this gap by combining FASHN-VTON for immediate 2D visualization with advanced 3D reconstruction models in a scalable web architecture suitable for real-world e-commerce use \cite{b11}.

\section{System Design and Methodology}

\subsection{Overall Workflow}
Stylmorph employs a bifurcated parallel processing pipeline designed to balance speed and depth of visualization. Upon uploading a full-body image, the system simultaneously dispatches data to two distinct pipelines. The primary pipeline ports the image to a remote Google Colaboratory environment via a secure tunneling connection, where the FASHN-VTON framework \cite{b23} performs garment warping and synthesis to produce a photorealistic 2D try-on result. The secondary pipeline runs PIFuHD \cite{b1} locally within a Python subprocess to generate a dense 3D `.obj` mesh encoding the user's full body topology. Both outputs are uploaded to Cloudinary on completion and their URLs are stored in MongoDB, linked to the user's profile.

\subsection{System Architecture}

The Stylmorph system follows a modular client-server architecture. The React.js frontend (built with Vite and Tailwind CSS) communicates over REST with an Express.js backend. The backend manages authentication via JWT, handles file uploads using Multer, and orchestrates ML execution via Node.js \texttt{child\_process} calls. For the 2D pipeline, Colab execution is triggered programmatically; for the 3D pipeline, a local Python environment runs PIFuHD inference. Output assets are offloaded to Cloudinary to prevent local disk bottlenecks.

\begin{figure}[h]
    \centering
    \includegraphics[width=0.48\textwidth]{Architecture_Diagram(fixed).png}
    \caption{Architecture diagram showing the decoupled 2D (Colab) and 3D (local PIFuHD) pipelines with CDN-backed storage.}
    \label{fig:architecture}
\end{figure}

\subsection{Modules}

\textbf{Frontend Module:} Developed with React.js, Vite, and Tailwind CSS. Integrates \texttt{@react-three/fiber} and Three.js \cite{b8} to provide real-time, hardware-accelerated 360-degree WebGL rendering of 3D avatars directly in the browser without any plugins.

\textbf{Backend Module:} Built on Node.js and Express.js. Handles RESTful API routing, JWT-based user authentication, image upload management via Multer, and subprocess spawning for ML script execution.

\textbf{Machine Learning Module:} Implemented in Python using PyTorch. The FASHN-VTON inference framework \cite{b23}, running on an external Colab GPU environment due to local infrastructure constraints, handles high-fidelity 2D garment transfer. The PIFuHD model \cite{b1}, run locally via its serialized checkpoint (\texttt{pifuhd.pt}), performs monocular implicit 3D reconstruction and outputs \texttt{.obj} mesh files.

\textbf{Database and Storage Module:} MongoDB with Mongoose ODM stores user profiles, measurement data, try-on histories, and cloud asset URLs. Cloudinary hosts and serves all generated 2D images and 3D mesh files via its global CDN.

\begin{table}[htbp]
\centering
\caption{Avatar Data Schema Structure}
\begin{adjustbox}{width=\linewidth}
\begin{tabular}{|l|p{5.5cm}|}
\hline
\textbf{Field} & \textbf{Description} \\
\hline
UserId & Unique identifier linking the avatar to a user account \\
\hline
ImageUrl & Cloudinary URL of the user's uploaded 2D source photo \\
\hline
VtonResults & Array of 2D try-on images from FASHN-VTON \cite{b23} with garment metadata \\
\hline
PifuhdUrl & Cloudinary URL of the generated PIFuHD \texttt{.obj} file \cite{b1} \\
\hline
Wearables & Array of 3D clothing items (Name, URL, Thumbnail) applied to avatar \\
\hline
SavedSets & Categorized saved outfit configurations \\
\hline
Measurements & Optional structured inputs (Height, Chest, Waist, etc.) \\
\hline
\end{tabular}
\end{adjustbox}
\label{tab:schema}
\end{table}

\section{Experimental Setup}

The Stylmorph experimental environment was configured as follows:

\textbf{Development Stack:} Frontend built using React 19, Vite 7, and Tailwind CSS 4. Backend powered by Node.js with Express.js 5 and MongoDB via Mongoose 8. Three.js 0.180 is used for 3D rendering via \texttt{@react-three/fiber}.

\textbf{ML Runtime:} The PIFuHD prediction pipeline runs locally using a pre-trained checkpoint (\texttt{pifuhd.pt}) within a Python virtual environment with PyTorch, OpenCV, Trimesh, and PyRender. No local model training was performed; the pipeline operates in inference-only mode. Due to hardware constraints, the FASHN-VTON 2D try-on pipeline was deployed on Google Colaboratory with a GPU runtime (T4), accessed via secure API tunneling from the Node.js backend.

\textbf{Cloud Infrastructure:} Cloudinary is used for storing and serving all generated 2D images and 3D \texttt{.obj} files. MongoDB Atlas serves as the cloud database for user data and asset references.

\section{Results and Discussion}

\subsection{Performance Analysis}

The parallel pipeline architecture successfully decouples the slow 2D Colab inference from the local 3D reconstruction, ensuring both run concurrently. The FASHN-VTON pipeline running over Google Colab produced high-fidelity try-on images preserving fine garment textures, logos, and folds across varied body poses \cite{b23}. The PIFuHD pipeline effectively extracted detailed 3D body topology from single monocular images, capturing clothing wrinkles and hair details in the output mesh \cite{b5}. Three.js rendered the resulting meshes smoothly in the browser at stable frame rates, even for complex meshes exceeding 50,000 polygon faces.

\subsection{Execution Latency}

Table \ref{tab:comparison} presents the measured execution latencies across all pipeline stages during system testing.

\begin{table}[htbp]
\centering
\caption{Pipeline Execution Latency and Output Metrics}
\begin{adjustbox}{width=\linewidth}
\begin{tabular}{|l|c|c|p{2.8cm}|}
\hline
\textbf{Pipeline Stage} & \textbf{Latency} & \textbf{Output} & \textbf{User Impact} \\
\hline
FASHN-VTON (Colab) & 7--8 min & JPEG Image & 2D garment visualization \\
\hline
PIFuHD 3D Reconstruction & 3--4 min & \texttt{.obj} Mesh & Full 360° avatar fit view \\
\hline
Three.js WebGL Rendering & $<$2 sec & Canvas (WebGL) & Interactive 3D exploration \\
\hline
\end{tabular}
\end{adjustbox}
\label{tab:comparison}
\end{table}

The 2D FASHN-VTON pipeline requires 7--8 minutes owing to the remote Colab execution overhead and network transfer latency. The 3D PIFuHD reconstruction completes locally in 3--4 minutes. Both pipelines run concurrently, so the user can inspect the 3D avatar while the 2D try-on is still processing, minimizing perceived wait time.

\subsection{Qualitative Results}

\begin{figure}[h]
    \centering
    \includegraphics[width=0.48\textwidth]{fashn_vton_output.png}
    \caption{FASHN-VTON 2D virtual try-on output. Left: uploaded person image (428$\times$612). Centre: selected garment image (512$\times$683). Right: synthesized try-on result (576$\times$823), preserving garment texture, logo, and overall fit.}
    \label{fig:fashn_output}
\end{figure}

\begin{figure}[h]
    \centering
    \includegraphics[width=0.45\textwidth]{pifu.png}
    \caption{PIFuHD 3D reconstruction output showing detailed surface geometry and clothing topology from a single monocular input image.}
    \label{fig:pifu}
\end{figure}

Figures \ref{fig:fashn_output} and \ref{fig:pifu} collectively demonstrate the two-stage output of the Stylmorph pipeline. The FASHN-VTON result confirms high-fidelity texture preservation and accurate garment morphing onto the target person's pose. The PIFuHD mesh validates the system's ability to produce spatially accurate 3D avatars suitable for comprehensive garment fitting analysis.


\subsection{Comparative Analysis}
Compared to traditional platforms relying solely on 2D image overlays \cite{b24}\cite{b25}, Stylmorph provides significant depth-perception benefits via 3D volumetric rotation. Purely 3D systems \cite{b1} require full processing time before any feedback is delivered. Stylmorph's concurrently running pipelines mitigate this by allowing the user to interact with the faster-completing 3D avatar while the detailed 2D try-on finalizes, offering the best characteristics of both approaches.

\subsection{Limitations}
The current system's primary limitation is the significant processing latency introduced by routing FASHN-VTON inference through a remote Colab environment. Input image quality—lighting, pose, and occlusions—also directly affects both 2D synthesis and 3D reconstruction fidelity. Future versions will prioritize migrating the 2D inference pipeline to a dedicated GPU server to dramatically reduce latency.

\section{Conclusion and Future Scope}

Stylmorph successfully demonstrates the feasibility of integrating state-of-the-art 2D virtual try-on (FASHN-VTON \cite{b23}) with explicit monocular 3D reconstruction (PIFuHD \cite{b1}) within a scalable, consumer-accessible web application. The MERN-stack architecture, augmented by Three.js WebGL rendering and Cloudinary CDN delivery, proves that high-fidelity fashion visualization is viable for real-world e-commerce deployment. By providing both immediate 2D garment previews and detailed 3D avatars, Stylmorph offers a significant advancement over existing single-modality approaches.

Future work will focus on three key areas: (1) migrating the FASHN-VTON pipeline to a dedicated GPU inference server to reduce latency from minutes to seconds; (2) projecting FASHN-VTON 2D texture maps directly onto PIFuHD 3D mesh surfaces \cite{b19} to produce fully textured photorealistic avatars; and (3) implementing physics-based cloth simulation within the Three.js ecosystem \cite{b20} to accurately model fabric draping, tension, and collision against the generated avatar mesh.

\section*{Acknowledgment}
The authors would like to thank the Department of Computer Science and Engineering, College of Engineering, Cherthala for providing the necessary infrastructure and support for this research work.

\begin{thebibliography}{00}
\bibitem{b1} Saito, S., Huang, Z., Natsume, R., et al. (2020). ``PIFuHD: Multi-Level Pixel-Aligned Implicit Function for High-Resolution 3D Human Digitization.'' IEEE/CVF CVPR.

\bibitem{b3} Loper, M., Mahmood, N., Romero, J., Pons-Moll, G., \& Black, M. J. (2015). ``SMPL: A Skinned Multi-Person Linear Model.'' ACM TOG, 34(6), 1--16.

\bibitem{b4} Dong, Z., et al. (2023). ``AG3D: Learning to Generate 3D Avatars from 2D Images.'' IEEE TPAMI, 45(8), 11234--11248.

\bibitem{b5} Ma, Q., et al. (2022). ``CAPE: Clothed Auto-Personalized Avatars.'' ACM TOG, 41(4), 1--15.

\bibitem{b7} Batool, R., \& Mou, J. (2023). ``Virtual Try-On Technology: A Comprehensive Review and Research Agenda.'' IJHCI, 39(5), 1023--1045.

\bibitem{b8} Hosen, M. S., et al. (2023). ``Mastering 3D Modeling in Blender.'' Int. J. Computer Graphics, 15(2), 45--62.

\bibitem{b11} Jaleel, M. U. (2023). ``Impact of 3D Product Displays on E-Commerce Consumer Behavior.'' IJECS, 14(3), 201--220.

\bibitem{b12} Yang, H., Yao, H., \& Zhou, S. (2021). ``Accurate 3D Human Body Measurement from Monocular Images Using Silhouette Analysis.'' Pattern Recognition Letters, 142, 18--25.

\bibitem{b14} Kanazawa, A., Black, M. J., Jacobs, D. W., \& Malik, J. (2018). ``End-to-End Recovery of Human Shape and Pose.'' IEEE CVPR, 7122--7131.

\bibitem{b15} Pavlakos, G., et al. (2018). ``Learning to Estimate 3D Human Pose and Shape from a Single Color Image.'' IEEE CVPR, 459--468.

\bibitem{b16} Bogo, F., et al. (2016). ``Keep It SMPL: Automatic Estimation of 3D Human Pose and Shape.'' ECCV, 561--578.

\bibitem{b17} Lassner, C., Pons-Moll, G., \& Gehler, P. (2017). ``A Generative Model of People in Clothing.'' ICCV, 853--862.

\bibitem{b18} Yang, Z., et al. (2020). ``Graphonomy: Universal Human Parsing via Graph Transfer Learning.'' IEEE CVPR, 7450--7459.

\bibitem{b19} Pons-Moll, G., et al. (2015). ``Dyna: A Model of Dynamic Human Shape in Motion.'' ACM TOG, 34(4), 120:1--120:14.

\bibitem{b20} Guan, P., et al. (2012). ``Drape: Dressing Any Person.'' ACM TOG, 31(4), 35:1--35:10.

\bibitem{b23} FASHN (2024). ``FASHN-VTON: High-Fidelity Virtual Try-On with Texture Preservation.'' FASHN Research Technical Report. Available: https://fashn.ai/research/vton

\bibitem{b24} Han, X., et al. (2018). ``VITON: An Image-based Virtual Try-on Network.'' IEEE CVPR, 7543--7552.

\bibitem{b25} Wang, B., et al. (2018). ``Toward Characteristic-Preserving Image-based Virtual Try-On Network.'' ECCV, 589--604.

\end{thebibliography}

\end{document}
